% ch1.tex -- Module 7, Chapter 1: The electromagnetic spectrum
% (Maxwell, the speed of light and spectroscopy). Dot points M7.IQ1.01-M7.IQ1.06.
% Figures: figures/ch1_*.py (matplotlib) and inline TikZ.
\chapter[The Electromagnetic Spectrum]{The Electromagnetic Spectrum: Maxwell, the Speed of Light and Spectroscopy}\label{ch1:chap}
\covers{M7.IQ1.01, M7.IQ1.02, M7.IQ1.03, M7.IQ1.04, M7.IQ1.05, M7.IQ1.06}

% Rainbow shading for the small TikZ spectrum strips in this chapter
% (violet on the left, red on the right; the middle 50bp maps onto the path).
\pgfdeclarehorizontalshading{ch1rainbow}{100bp}{%
  color(0bp)=(violet); color(25bp)=(violet); color(31bp)=(blue);
  color(38bp)=(cyan); color(45bp)=(green!85!black); color(56bp)=(yellow);
  color(63bp)=(orange); color(70bp)=(red); color(75bp)=(red); color(100bp)=(red)}

\inquiry{What is light?}

\begin{goals}
  \item explain how Maxwell unified electricity and magnetism, and how his theory predicted
        electromagnetic waves and their speed
  \item describe how electromagnetic waves are made and how they travel; use $c = f\lambda$
  \item outline how the speed of light was measured, and how it now defines the metre
  \item describe continuous, emission and absorption spectra, and use them to identify elements
  \item explain what a star's spectrum tells us about its temperature, motion, rotation,
        density and chemical composition
\end{goals}

%% ==================================================================
\section{Maxwell and electromagnetic waves}\label{ch1:sec:maxwell}
\covers{M7.IQ1.01, M7.IQ1.02}

\begin{bigidea}
Maxwell joined electricity and magnetism into one theory of four laws. The laws predicted waves
of electric and magnetic fields, made by accelerating charges and travelling at
\qty{3.00e8}{m.s^{-1}}. This is the speed of light, so light is an electromagnetic wave.
\end{bigidea}

\begin{keywords}
  \kw{Unification}{showing that two things that look different are parts of one thing.}
  \kw{Electromagnetic (EM) wave}{a travelling wave of electric and magnetic fields that change
    together.}
  \kw{Accelerating charge}{a charge that speeds up, slows down or changes direction.}
  \kw{Transverse wave}{a wave whose vibrations are at \ang{90} to its direction of travel.}
  \kw{In phase}{reaching their peaks at the same time and the same place.}
\end{keywords}

\begin{refresher}{Fields and waves (Module 6 and Year 11)}
\begin{itemize}
  \item An \term{electric field} is a region where a charge feels a force. A \term{magnetic
        field} is a region where a magnet or a moving charge feels a force.
  \item A current makes a magnetic field. A changing magnetic flux (the amount of magnetic
        field passing through a loop) makes an emf, which is a voltage (Faraday's law).
  \item For a wave: wavelength $\lambda$ (\si{m}), frequency $f$ (\si{Hz}), period $T$
        (\si{s}), with $f = \dfrac{1}{T}$ and $v = f\lambda$. Both formulas are on the sheet.
\end{itemize}
\end{refresher}

Before the 1860s, scientists knew some links between electricity and magnetism. A current makes
a magnetic field (Oersted, 1820). A changing magnetic field makes a current (Faraday, 1831). But
there was no single theory, and light seemed to be a separate topic.

Most scientists thought light was a wave in the \term{luminiferous aether}. This was an
invisible material that they believed filled all space.

Between 1861 and 1865, James Clerk Maxwell wrote all of electricity and magnetism as one set of
laws. (Maxwell's own papers used about 20 equations. Oliver Heaviside later wrote them as the
four laws we use today, \cref{ch1:tab:laws}.) You need to know what each law says, not the
equations.

\begin{table}
  \caption{Maxwell's four laws in plain words.}\label{ch1:tab:laws}
  \centering
  \begin{tabular}{@{}L{44mm}L{100mm}@{}}
    \toprule
    \colhead{Law} & \colhead{What it says}\\
    \midrule
    1.\ Gauss's law (electric) & Charges make electric fields. Field lines start on $+$ charges
      and end on $-$ charges.\\
    2.\ Gauss's law for magnetism & Magnetic field lines always form closed loops. A single N or
      S pole on its own does not exist.\\
    3.\ Faraday's law & A changing magnetic field makes an electric field, even in empty space.\\
    4.\ Ampère--Maxwell law & A magnetic field is made by a current \emph{and} by a changing
      electric field.\\
    \bottomrule
  \end{tabular}
\end{table}

The second half of law 4 was Maxwell's own new idea. With it, laws 3 and 4 become a matching
pair. A changing magnetic field makes an electric field (law 3). A changing electric field
makes a magnetic field (law 4). Electric and magnetic fields are two parts of one \term{electromagnetic field}. This
is Maxwell's \term{unification} of electricity and magnetism.

\subsection{Maxwell's two predictions}

\begin{sidebyside}[0.50]
  \centering
  \begin{tikzpicture}
    \node[concept=23mm] (E) at (0,0) {changing\\electric field};
    \node[concept=23mm] (B) at (3.9,0) {changing\\magnetic field};
    \draw[link] (E.north) to[bend left=35] node[linklabel, above] {makes (law 4)} (B.north);
    \draw[link] (B.south) to[bend left=35] node[linklabel, below] {makes (law 3)} (E.south);
  \end{tikzpicture}
  \captionof{figure}{Each changing field makes the other.}\label{ch1:fig:chain}
\nextside
  \textbf{1. EM waves.} A changing electric field makes a changing magnetic field. That makes a
  changing electric field a little further on, and so on (\cref{ch1:fig:chain}). The pair keeps
  itself going and travels away. It needs no wire and no charges. Nothing limits the frequency,
  so Maxwell predicted many kinds of EM waves, most of them invisible.
\end{sidebyside}

\textbf{2. Their speed.} The laws give the speed of the waves from two constants. The
\term{permeability of free space} ($\mu_0$) sets how strong magnetic fields are in a vacuum.
The \term{permittivity of free space} ($\varepsilon_0$) sets how strong electric fields are in
a vacuum. ``Free space'' means a vacuum: space with no air or other matter. Both constants are
measured in ordinary laboratory experiments with currents and charges, not with light.

In 1862, Maxwell calculated a speed of about \qty{3.1e8}{m.s^{-1}}. Fizeau had measured the
speed of light in 1849 and also got about \qty{3.1e8}{m.s^{-1}}. Maxwell concluded that light is
an EM wave. (Maxwell still believed in the aether. Later experiments showed that light needs no
medium, \cref{ch4:chap}.)

\begin{toolkit}{Powers of ten, square roots and percentages}
\begin{itemize}
  \item Prefixes: n $=\times10^{-9}$, \ $\upmu$ $=\times10^{-6}$, \ m $=\times10^{-3}$,
        \ k $=\times10^{3}$, \ M $=\times10^{6}$, \ G $=\times10^{9}$. So
        $\qty{525}{nm} = \qty{5.25e-7}{m}$. Time: \qty{1}{min} $=$ \qty{60}{s}.
  \item Multiply powers of ten by adding the powers: $10^{-7}\times10^{-12} = 10^{-19}$. A
        square root halves the power: $\sqrt{10^{-18}} = 10^{-9}$.
  \item If the power is odd, move the decimal point first: $1.11\times10^{-17} =
        11.1\times10^{-18}$, so $\sqrt{11.1\times10^{-18}} \approx 3.33\times10^{-9}$.
  \item For $\dfrac{1}{\sqrt{\mu_0\varepsilon_0}}$, type $1 \div \sqrt{(\mu_0 \times
        \varepsilon_0)}$. Use the $\times10^{x}$ key for powers of ten.
  \item Percentage difference $= \dfrac{\text{difference}}{\text{accepted value}}\times100\,\%$.
\end{itemize}
\end{toolkit}

\begin{toolkit}{Rearranging a formula}
\begin{itemize}
  \item To get one symbol on its own, do the same thing to both sides.
  \item Example: $c = f\lambda$. Divide both sides by $f$: $\dfrac{c}{f} = \lambda$. So
        $\lambda = \dfrac{c}{f}$.
  \item Example: $\lambda_\mathrm{max} = \dfrac{b}{T}$. Multiply both sides by $T$:
        $\lambda_\mathrm{max}T = b$. Divide both sides by $\lambda_\mathrm{max}$:
        $T = \dfrac{b}{\lambda_\mathrm{max}}$.
  \item Check: put numbers into the old and the new formula. Both should give the same answer.
\end{itemize}
\end{toolkit}

\begin{formulacard}[notonsheet]{Speed of electromagnetic waves (Maxwell)}
  \cardformula{c = \frac{1}{\sqrt{\mu_0\varepsilon_0}}}
  \begin{symbols}
    \sym{$c$}{speed of EM waves in a vacuum}{\si{m.s^{-1}}}
    \sym{$\mu_0$}{permeability of free space, $4\pi\times10^{-7}$ (data sheet)}{\si{N.A^{-2}}}
    \sym{$\varepsilon_0$}{permittivity of free space, \num{8.854e-12} (data sheet)}%
      {\si{A^2.s^4.kg^{-1}.m^{-3}}}
  \end{symbols}
  \inwords{the speed of light comes from one electric constant and one magnetic constant.}
  \whentouse{to calculate Maxwell's predicted speed, or to show where it comes from.}
  \whennottouse{for light inside glass or water, where light is slower.}
  \rearranged{$\mu_0\varepsilon_0 = \dfrac{1}{c^2}$ \qquad $\varepsilon_0 = \dfrac{1}{\mu_0 c^2}$}
  \tensecond{with $c = \qty{3.00e8}{m.s^{-1}}$:
    $\varepsilon_0 = \dfrac{1}{4\pi\times10^{-7}\times(3.00\times10^{8})^2}
    = \qty{8.84e-12}{A^2.s^4.kg^{-1}.m^{-3}}$. This is close to the data sheet's
    \num{8.854e-12}. It is a little lower because $c = \qty{3.00e8}{m.s^{-1}}$ is rounded.}
  \derive{for empty space, laws 3 and 4 combine (with university maths) into a wave whose speed
    $v$ obeys $v^2\mu_0\varepsilon_0 = 1$. Divide both sides by $\mu_0\varepsilon_0$:
    $v^2 = \dfrac{1}{\mu_0\varepsilon_0}$. Take the square root of both sides:
    $v = \dfrac{1}{\sqrt{\mu_0\varepsilon_0}}$. You only need the result. Unit check, with
    \qty{1}{N} $=$ \qty{1}{kg.m.s^{-2}}: $\mu_0\varepsilon_0$ has unit
    $(\unit{kg.m.s^{-2}.A^{-2}})(\unit{A^2.s^4.kg^{-1}.m^{-3}}) = \unit{s^2.m^{-2}}$. So
    $\dfrac{1}{\sqrt{\mu_0\varepsilon_0}}$ has unit \si{m.s^{-1}}, a speed.}
  \trap{forgetting the square root gives \num{8.99e16}, which is $c^2$, not $c$.}
\end{formulacard}

\begin{workedexample}{The missing square root}\label{ch1:we:maxwell}
A student uses the data sheet values of $\mu_0$ and $\varepsilon_0$. She types
$1 \div (\mu_0\times\varepsilon_0)$ and forgets the square root. (a) What number does she get?
How can she tell that it is not a speed? (b) Calculate the correct speed of EM waves. Compare it
with $c$ on the data sheet.
\given{$\mu_0 = 4\pi\times10^{-7}$ \si{N.A^{-2}}, \
  $\varepsilon_0 = \qty{8.854e-12}{A^2.s^4.kg^{-1}.m^{-3}}$}
\find{(a) her number; (b) the predicted speed $c$}
\formula{$c = \dfrac{1}{\sqrt{\mu_0\varepsilon_0}}$}
\substitute{multiply first:\\
  \eqlines{\mu_0\varepsilon_0 &= (4\pi\times10^{-7})\times(8.854\times10^{-12})
    = \qty{1.1126e-17}{s^2.m^{-2}}}\\
  (a) her keys, with no square root:\\
  \eqlines{\frac{1}{\mu_0\varepsilon_0} &= \frac{1}{\qty{1.1126e-17}{s^2.m^{-2}}}
    = \qty{8.988e16}{m^2.s^{-2}}}\\
  (b) take the square root, then divide 1 by it:\\
  \eqlines{\sqrt{\mu_0\varepsilon_0} &= \sqrt{\qty{1.1126e-17}{s^2.m^{-2}}}
    = \qty{3.3356e-9}{s.m^{-1}}\\
  c &= \frac{1}{\qty{3.3356e-9}{s.m^{-1}}} = \qty{2.998e8}{m.s^{-1}}}}
\answer{(a) \final{\qty{8.99e16}{m^2.s^{-2}}}. The unit is \si{m^2.s^{-2}}, not
  \si{m.s^{-1}}, so it is not a speed. It is $c^2$. (b) \final{$c = \qty{2.998e8}{m.s^{-1}}$}.
  This agrees with \qty{3.00e8}{m.s^{-1}} to 3 significant figures.}
\checkstep{the difference is $3.00 - 2.998 = 0.002$ (in units of \qty{e8}{m.s^{-1}}). As a
  percentage: $\dfrac{0.002}{3.00}\times100\,\% = 0.07\,\%$. Electric and magnetic constants
  give the measured speed of light: this is Maxwell's evidence that light is an EM wave.}
\end{workedexample}

\subsection{How EM waves are made}

A stone dropped in a pond makes ripples. An accelerating charge does the same to its field:
\begin{itemize}
  \item A charge at rest, or moving at a \emph{constant} velocity, has steady fields. It sends
        out no wave.
  \item A charge that \emph{accelerates} disturbs its field. The disturbance spreads out at $c$.
  \item An oscillating charge (one that moves back and forth, again and again) accelerates all
        the time. It sends out a continuous EM wave with the \emph{same frequency} as its
        oscillation.
\end{itemize}

In a radio transmitter, an AC (alternating current) supply pushes electrons back and forth in
the \term{antenna}, a metal rod that sends out or picks up radio waves (\cref{ch1:fig:antenna}).
When the wave reaches a receiving antenna, its electric field pushes the electrons there back
and forth at the same frequency. This makes a small AC current.

\begin{figure}
  \centering
  \begin{tikzpicture}[x=1cm, y=1cm]
    \draw[line width=2.2pt, hscink] (0,0.2) -- (0,1.3);
    \draw[line width=2.2pt, hscink] (0,-0.2) -- (0,-1.3);
    \draw[line width=0.7pt] (0,0.2) -| (-0.9,0.28);
    \draw[line width=0.7pt] (0,-0.2) -| (-0.9,-0.28);
    \draw[line width=0.7pt, fill=white] (-0.9,0) circle (0.28);
    \draw[line width=0.7pt] (-1.06,0) sin (-0.98,0.09) cos (-0.9,0) sin (-0.82,-0.09) cos (-0.74,0);
    \node[diagramlabel, align=center, anchor=east] at (-1.25,0) {AC\\supply};
    \draw[<->, line width=1.1pt, hscacc] (0.26,0.8) -- (0.26,-0.8);
    \foreach \r in {1.3,1.95,2.6,3.25}
      \draw[hscfield!75, line width=0.9pt] (-24:\r) arc (-24:24:\r);
    \draw[->, line width=0.9pt, hscink] (3.6,0) -- (4.5,0) node[midway, above, diagramlabel] {at $c$};
    \draw[line width=2.2pt, hscink] (5.4,0.18) -- (5.4,1.15);
    \draw[line width=2.2pt, hscink] (5.4,-0.18) -- (5.4,-1.15);
    \draw[line width=0.7pt] (5.4,0.18) -| (6.2,0.28);
    \draw[line width=0.7pt] (5.4,-0.18) -| (6.2,-0.28);
    \draw[line width=0.7pt, fill=white] (6.2,0) circle (0.28) node[font=\small\bfseries] {A};
    \node[diagramlabel, anchor=west] at (6.55,0) {ammeter};
    \draw[<->, line width=1.1pt, hscacc] (5.14,0.7) -- (5.14,-0.7);
    \node[diagramlabel, anchor=north, align=center] at (0,-1.4) {transmitter: charges oscillate};
    \node[diagramlabel, anchor=north, align=center] at (5.7,-1.4) {receiver: charges oscillate\\at the same frequency};
  \end{tikzpicture}
  \caption{A radio transmitter and receiver. The arcs are crests of the wave. The ammeter
    detects the small AC current in the receiver.}
  \label{ch1:fig:antenna}
\end{figure}

\subsection{How EM waves travel}

\begin{sidebyside}[0.56]
  \centering
  \begin{tikzpicture}[x={(0.82cm,0cm)}, y={(0cm,0.82cm)}, z={(-0.24cm,-0.30cm)}]
    \def\lam{3.0}
    \fill[hscseriesA!7] (0,-1.25,0) -- (6,-1.25,0) -- (6,1.25,0) -- (0,1.25,0) -- cycle;
    \fill[hscforce!15] (0,0,-1.25) -- (6,0,-1.25) -- (6,0,1.25) -- (0,0,1.25) -- cycle;
    \foreach \x in {0.125,0.25,...,6.01}
      \draw[hscseriesA!55, line width=0.35pt] (\x,0,0) -- (\x,{1.25*sin(360*\x/\lam)},0);
    \draw[hscseriesA, line width=1.3pt, domain=0:6, samples=121, smooth]
      plot (\x,{1.25*sin(360*\x/\lam)},0);
    \foreach \x in {0.125,0.25,...,6.01}
      \draw[hscforce!55, line width=0.35pt] (\x,0,0) -- (\x,0,{1.25*sin(360*\x/\lam)});
    \draw[hscforce, line width=1.3pt, domain=0:6, samples=121, smooth]
      plot (\x,0,{1.25*sin(360*\x/\lam)});
    \draw[axisline] (0,0,0) -- (6.9,0,0) node[below, diagramlabel, align=center]
      {direction\\of travel};
    \draw[axisline] (0,0,0) -- (0,1.7,0);
    \draw[axisline] (0,0,0) -- (0,0,2.0);
    \node[diagramlabel, text=hscseriesA, anchor=west] at (4.05,1.05,0) {$E$};
    \node[diagramlabel, text=hscforce, anchor=north east] at (0.65,0,1.3) {$B$};
    \draw[<->, line width=0.6pt, hscink] (0.75,1.62,0) -- (3.75,1.62,0)
      node[midway, above, diagramlabel] {$\lambda$};
    \draw[guide] (0.75,1.3,0) -- (0.75,1.68,0);
    \draw[guide] (3.75,1.3,0) -- (3.75,1.68,0);
    % one pair of crests: E and B are largest at the same place (in phase)
    \draw[hscink, line width=0.6pt, dashed] (3.75,1.25,0) -- (3.75,0,0) -- (3.75,0,1.25);
    \fill[hscink] (3.75,1.25,0) circle (1.6pt);
    \fill[hscink] (3.75,0,1.25) circle (1.6pt);
    \draw[hscforce, line width=0.5pt] (3.75,0,1.25) -- ++(0cm,-0.66cm)
      node[below, inner sep=1pt, diagramlabel, text=hscforce] {$B$ largest};
  \end{tikzpicture}
  \captionof{figure}{An EM wave at one instant. $E$ (blue) is in the vertical plane; $B$ (red)
    is in the horizontal plane, drawn coming out towards you. The black dots: both are largest
    at the same place.}%
  \label{ch1:fig:emwave}
\nextside
  \Cref{ch1:fig:emwave} shows that:
  \begin{itemize}
    \item $E$ and $B$ oscillate with the same frequency, and \term{in phase}.
    \item $E$ and $B$ are at \ang{90} to each other and to the direction of travel: a
          \term{transverse wave}.
    \item It can travel through a vacuum (no medium is needed). In a vacuum it travels at $c$.
    \item It carries energy.
  \end{itemize}
\end{sidebyside}

All of these follow from Maxwell's laws. Each changing field makes the other (laws 3 and 4), so
the wave can travel through empty space. The laws also predict that $E$, $B$ and the direction
of travel are all at \ang{90} to each other. They also give the speed
$\dfrac{1}{\sqrt{\mu_0\varepsilon_0}} = c$. (Maxwell himself still pictured the wave in the
aether. Experiments later showed that no medium is needed, \cref{ch4:chap}.)

In 1887, Heinrich Hertz made radio waves with an oscillating spark. He detected them across the
room with a loop of wire that had a tiny gap. Small sparks jumped across the gap when the waves
arrived. The waves reflected and refracted, like light. They were also polarised: their electric
field vibrated in one direction only (\cref{ch2:chap}). Their speed was close to $c$. This
confirmed Maxwell's theory.

\subsection{The electromagnetic spectrum}

All EM waves are the same kind of wave. They differ only in frequency, and so in wavelength.
The \term{electromagnetic spectrum} is the whole range (\cref{ch1:fig:emspectrum}). Visible
light is a tiny part, from about \qty{400}{nm} (violet) to \qty{700}{nm} (red). Every source
involves charges: accelerating in an antenna, or changing energy inside atoms and nuclei.

\begin{figure}
  \centering
  \includegraphics{ch1_emspectrum.pdf}
  \caption{The electromagnetic spectrum, with a typical source of each band. Band edges are
    approximate.}\label{ch1:fig:emspectrum}
\end{figure}

\begin{formulacard}[combine={$v = f\lambda$ with $v = c$}]{Wave equation for light}
  \cardformula{c = f\lambda}
  \begin{symbols}
    \sym{$c$}{speed of light in a vacuum, \qty{3.00e8}{m.s^{-1}}}{\si{m.s^{-1}}}
    \sym{$f$}{frequency}{\si{Hz}}
    \sym{$\lambda$}{wavelength}{\si{m}}
  \end{symbols}
  \inwords{speed = (number of waves each second) $\times$ (length of one wave).}
  \whentouse{any EM wave, from radio to gamma rays, in a vacuum or in air.}
  \whennottouse{light inside glass or water, where it is slower (use $v = f\lambda$).}
  \rearranged{$f = \dfrac{c}{\lambda}$ \qquad $\lambda = \dfrac{c}{f}$ \qquad $T = \dfrac{1}{f}$}
  \tensecond{a Wi-Fi signal has $f = \qty{2.40}{GHz} = \qty{2.40e9}{Hz}$, so
    $\lambda = \dfrac{3.00\times10^{8}}{2.40\times10^{9}} = \qty{0.125}{m}$ (microwaves).}
  \derive{the sheet gives $v = f\lambda$ for every wave. For EM waves in a vacuum, $v = c$.}
  \trap{change \si{nm} to \si{m} ($\times10^{-9}$), and \si{MHz} or \si{GHz} to \si{Hz},
    before you substitute.}
\end{formulacard}

\begin{workedexample}{From the period to the colour}\label{ch1:we:period}
At one point, the electric field of an EM wave makes one full oscillation in
\qty{1.75e-15}{s}. (a) Calculate the frequency and the wavelength. (b) Identify the radiation.
(c) Describe the magnetic field of the wave at this point.
\given{$T = \qty{1.75e-15}{s}$, \ $c = \qty{3.00e8}{m.s^{-1}}$}
\find{$f$ and $\lambda$, then the type of radiation}
\formula{$f = \dfrac{1}{T}$ \quad then \quad $\lambda = \dfrac{c}{f}$}
\substitute{\eqlines{f &= \frac{1}{\qty{1.75e-15}{s}} = \qty{5.714e14}{Hz}
    \quad\text{(keep the extra digits)}\\
  \lambda &= \frac{\qty{3.00e8}{m.s^{-1}}}{\qty{5.714e14}{Hz}} = \qty{5.25e-7}{m}
    = \qty{525}{nm}}}
\answer{(a) \final{$f = \qty{5.71e14}{Hz}$}, \final{$\lambda = \qty{525}{nm}$}. (b) Visible
  light (green). (c) $B$ oscillates at the same frequency, in phase with $E$, at \ang{90} to
  $E$ and to the direction of travel.}
\checkstep{\qty{525}{nm} is between 400 and \qty{700}{nm}, so the light is visible \cmark.}
\end{workedexample}

\begin{mistakes}
\begin{itemize}
  \item \xmark\ ``Maxwell measured $c$'' or ``Maxwell discovered radio waves.'' \cmark\ He
        \emph{predicted} both. Experimenters such as Fizeau and Foucault measured $c$; Hertz
        made radio waves.
  \item \xmark\ Law 4 is only ``a current makes a magnetic field''. \cmark\ A changing
        electric field does too. Without this part there are no EM waves.
  \item \xmark\ ``A charge moving at a steady speed sends out EM waves.'' \cmark\ Only
        accelerating charges do.
  \item \xmark\ ``X-rays travel faster than radio waves.'' \cmark\ All EM waves travel at $c$
        in a vacuum.
\end{itemize}
\end{mistakes}

\begin{examtip}[breakable]
\textbf{``Analyse Maxwell's contribution'' (5--6 marks).} Show how each idea leads to the next:
\begin{enumerate}
  \item Before Maxwell: some links were known (Oersted, Faraday), but there was no single
        theory. Light was a separate topic.
  \item He unified electricity and magnetism in four laws. His new idea: a changing $E$ field
        makes a $B$ field.
  \item So he predicted self-sustaining EM waves of many frequencies.
  \item Their predicted speed, $\dfrac{1}{\sqrt{\mu_0\varepsilon_0}} = \qty{3.00e8}{m.s^{-1}}$,
        equals the measured $c$. So light is an EM wave.
  \item Hertz later made and detected radio waves. This confirmed the theory. (It was Hertz's
        work, not Maxwell's: a common multiple-choice trap.)
\end{enumerate}
\textbf{``Describe the production and propagation of EM waves'' (3--4 marks).} Accelerating
charges $\to$ changing $E$ and $B$ that make each other $\to$ in phase, at \ang{90} to each other
and to the direction of travel $\to$ speed $c$, no medium. In a sketch, label $E$, $B$,
$\lambda$ and the direction of travel.
\end{examtip}

\begin{quickcheck}
  \item State the part of the Ampère--Maxwell law that was Maxwell's own idea. Explain why it was
        needed to predict EM waves.\label{ch1:qc:addition}
  \item A student uses rounded values $\mu_0 = \qty{1.26e-6}{N.A^{-2}}$ and
        $\varepsilon_0 = \qty{8.85e-12}{A^2.s^4.kg^{-1}.m^{-3}}$. Calculate the speed she
        predicts for EM waves.
  \item Describe how a radio transmitter produces an EM wave. (3 marks)
  \item An EM wave travels east. At one instant its electric field points straight up. In which
        directions can its magnetic field point?
  \item An AM radio station broadcasts at \qty{702}{kHz}. Calculate the wavelength.
\end{quickcheck}

%% ==================================================================
\section{Measuring the speed of light}\label{ch1:sec:speed}
\covers{M7.IQ1.03}

\begin{bigidea}
Light crosses \qty{1}{km} in about 3 millionths of a second, so its speed is hard to measure.
Each new method measured a time or a distance better. Since 1983, the speed of light has an
exact, fixed value, and the metre is defined from it.
\end{bigidea}

\begin{keywords}
  \kw{Direct method}{measure a distance and a travel time, then calculate $c = d/t$.}
  \kw{Indirect method}{calculate $c$ from other measured quantities, for example $c = f\lambda$.}
  \kw{Atomic clock}{an extremely accurate clock that counts oscillations of radiation from atoms.}
\end{keywords}

\begin{refresher}{Working scientifically words}
\textbf{Validity}: the method really tests what it claims to, with other variables controlled.
\textbf{Reliability}: repeats give consistent results. \textbf{Accuracy}: how close a result is
to the true value. \textbf{Precision}: how close repeated results are to each other.
\textbf{Random error}: small, unpredictable changes between repeats; averaging many readings
reduces it.
\end{refresher}

The methods below measured a time and a distance better and better. \Cref{ch1:tab:history}, at
the end of this part, sums them up. The Key experiments (\cref{ch1:sec:experiments}) give the
details.

\subsection{How Rømer and Fizeau timed light}

\textbf{Rømer (1676).} Io is a moon of Jupiter. Every 1.77 days it moves into Jupiter's shadow
and disappears for a while. This is an \term{eclipse} of Io: the moon is hidden in a shadow.
As Earth moves to the far side of its orbit, the light from each eclipse has further to travel.
So we see the eclipses later and later (\cref{ch1:fig:romer}, page~\pageref{ch1:fig:romer}). The
extra distance is about the diameter of Earth's orbit, $2r$. The extra time is the delay.

\textbf{Fizeau (1849).} A \term{toothed wheel} is a wheel with equal teeth and gaps around its
edge, like a gear. Light goes out through a gap to a mirror far away and comes back. If the
wheel has turned a little, a tooth is now where the gap was. The tooth blocks the returning
light, and the observer sees darkness (\cref{ch1:fig:fizeau}, page~\pageref{ch1:fig:fizeau}). The
time for that tiny turn equals the travel time of the light.

\textbf{Rotating mirror (Foucault, Michelson).} Foucault and Michelson used a fast-spinning
mirror instead of a wheel. Light reflects off the spinning mirror, goes to a far mirror and
comes back. While the light travels, the spinning mirror turns a little. Foucault measured how
far the returning beam moved. Michelson spun his 8-sided mirror until it turned exactly
one-eighth of a turn during the trip. The spin rate gives the travel time, and so $c$.


\subsection{Today: the speed of light defines the metre}

\textbf{The 1972 laser method.} This is an indirect method. The frequency $f$ of a very stable
laser was counted against a caesium atomic clock. Its wavelength $\lambda$ was measured by
interference (\cref{ch2:chap}), against the krypton light that then defined the metre. Then
$c = f\lambda = \qty{299792456}{m.s^{-1}}$, uncertain by only about \qty{1}{m.s^{-1}}. Most of
that uncertainty came from the krypton definition of the metre itself. So in 1983 the definition
was turned around: $c$ was fixed, and the metre is now defined from $c$ and the second.

\begin{itemize}
  \item \textbf{The second} (since 1967) is \num{9192631770} oscillations of the radiation from
        one energy jump in a caesium-133 atom. An \term{atomic clock} counts them.
  \item \textbf{Since 1983}, $c$ is \emph{defined} as exactly \qty{299792458}{m.s^{-1}}. One
        \textbf{metre} is the distance light travels in a vacuum in $\dfrac{1}{299\,792\,458}$
        of a second. (From 1889 to 1960 the metre was the length of a metal bar kept near
        Paris. From 1960 to 1983 it was set by the wavelength of light from krypton atoms.)
  \item So we no longer measure $c$. We measure a \emph{length} by timing light with an atomic
        clock. Radar, GPS, laser rangefinders and the light-year all use this.
\end{itemize}
The data sheet value, \qty{3.00e8}{m.s^{-1}}, is the exact value to 3 significant figures.

\begin{table}
  \caption{Measuring the speed of light.}\label{ch1:tab:history}
  \centering\small
  \begin{tabular}{@{}L{24mm}L{39mm}L{31mm}L{42mm}@{}}
    \toprule
    \colhead{Who, when} & \colhead{Method} & \colhead{Result (\si{m.s^{-1}})} &
      \colhead{What limited it}\\
    \midrule
    Galileo, about 1638 & lanterns on two hills about \qty{1.5}{km} apart & no delay seen &
      reaction time is far longer than the travel time\\
    Rømer, 1676 & delays of the eclipses of Jupiter's moon Io & not instant; Huygens (1690):
      about \num{2.2e8} & eclipse timing; size of Earth's orbit\\
    Fizeau, 1849 & toothed wheel; mirror \qty{8.6}{km} away & about \num{3.1e8} & judging
      ``darkness'' by eye\\
    Foucault, 1862 & rotating mirror & \num{2.98e8} & measuring a tiny angle\\
    Michelson, 1926 & rotating 8-sided mirror; \qty{35}{km} each way & \num{2.99796e8} &
      air on the path (light is a little slower in air); uncertainty only
      \qty{4}{km.s^{-1}}\\
    Lasers, 1972 & measure $f$ and $\lambda$ of a laser; $c = f\lambda$ & \num{299792456}
      ($\pm1$) & the metre itself, then defined by krypton light (not precise enough)\\
    \bottomrule
  \end{tabular}
\end{table}

\begin{formulacard}[notonsheet]{Speed from distance and time}
  \cardformula{c = \frac{d}{t}}
  \begin{symbols}
    \sym{$c$}{speed of light}{\si{m.s^{-1}}}
    \sym{$d$}{distance travelled by the light}{\si{m}}
    \sym{$t$}{time taken}{\si{s}}
  \end{symbols}
  \inwords{speed = distance $\div$ time.}
  \whentouse{timing light over a known distance: Rømer, Fizeau, radar, GPS, laser ranging.}
  \whennottouse{you know the frequency and the wavelength: use $c = f\lambda$.}
  \rearranged{$d = ct$ \qquad $t = \dfrac{d}{c}$}
  \tensecond{a laser pulse to the Moon and back takes \qty{2.56}{s}. The Moon is
    $\dfrac{3.00\times10^{8}\times2.56}{2} = \qty{3.84e8}{m}$ away.}
  \derive{this is the definition of average speed. Special cases:
    there and back (distance $D$ each way), $d = 2D$;
    Rømer, $d = 2r$ (the diameter of Earth's orbit) and $t$ = the extra delay;
    Fizeau (wheel with $N$ teeth and $N$ gaps, turning $f$ times each second): the edge has
    $2N$ equal parts (teeth and gaps). Moving from a gap to the next tooth is one part, which is
    $\dfrac{1}{2N}$ of a turn. One full turn takes $\dfrac{1}{f}$ seconds. So
    $t = \dfrac{1}{2N}\times\dfrac{1}{f} = \dfrac{1}{2Nf}$.}
  \trap{a reflected signal goes there \emph{and back}: $d = 2D$. On the formula sheet, $s$ means
    displacement, so do not use $s$ for speed.}
\end{formulacard}

\begin{workedexample}{A Rømer-style calculation}\label{ch1:we:romer}
Io's eclipses are timed when Earth is closest to Jupiter. Six months later, Earth is on the far
side of its orbit, and the eclipses are \qty{16.4}{min} late. The radius of Earth's orbit is
\qty{1.50e11}{m} (see \cref{ch1:fig:romer}). (a) Estimate the speed of light. (b) Rømer's own
estimate of the delay was about \qty{22}{min}. What speed does this give? (This is a simplified
model. It ignores Jupiter's own motion. Also, Jupiter cannot be seen when it is behind the Sun,
so Rømer used eclipses timed in the months before and after.)
\given{$r = \qty{1.50e11}{m}$; change each delay from minutes to seconds:\\
  (a) $t = 16.4\times\qty{60}{s} = \qty{984}{s}$, \quad (b) $t = 22\times\qty{60}{s} =
  \qty{1320}{s}$}
\find{the speed of light $c$}
\formula{$c = \dfrac{d}{t}$, with $d = 2r$ (the diameter of Earth's orbit)}
\substitute{\eqlines{\text{(a)}\quad c &= \frac{2\times\qty{1.50e11}{m}}{\qty{984}{s}}
    = \qty{3.049e8}{m.s^{-1}}\\
  \text{(b)}\quad c &= \frac{2\times\qty{1.50e11}{m}}{\qty{1320}{s}}
    = \frac{\qty{3.00e11}{m}}{\qty{1320}{s}}\\
    &= \qty{2.27e8}{m.s^{-1}} \quad\text{(keep the extra digit for now)}}}
\answer{(a) \final{$c = \qty{3.05e8}{m.s^{-1}}$} \ (b) \final{$c = \qty{2.3e8}{m.s^{-1}}$}
  (2 significant figures, because \qty{22}{min} has 2). This
  is too low by $\dfrac{3.00 - 2.27}{3.00}\times100\,\% = 24\,\%$. Rømer's timing of the eclipses
  was a big source of error. (Huygens got about \qty{2.2e8}{m.s^{-1}}. He used Rømer's delay
  with the size of Earth's orbit known in his time, which was also poorly known.)}
\checkstep{(a) is within 2\,\% of \qty{3.00e8}{m.s^{-1}} \cmark. A longer delay over the same
  distance must give a smaller speed \cmark.}
\end{workedexample}

\begin{workedexample}{A Fizeau-style toothed wheel}\label{ch1:we:fizeau}
A wheel has 600 teeth and 600 equal gaps. The far mirror is \qty{5.20}{km} away. The returning
light is first blocked when the wheel turns 24.2 times per second (see \cref{ch1:fig:fizeau}).
Calculate the speed of light.
\given{$N = 600$, \ $D = \qty{5.20e3}{m}$, \ $f = \qty{24.2}{\rev.s^{-1}}$ (24.2 turns each
  second). Here $f$ is how fast the wheel turns, not a light frequency. A turn has no unit, so
  this is \qty{24.2}{s^{-1}}.}
\find{the speed of light $c$}
\formula{$t = \dfrac{1}{2Nf}$ (gap to next tooth, $\dfrac{1}{2N}$ of a turn), then
  $c = \dfrac{d}{t}$ with $d = 2D$}
\substitute{\eqlines{t &= \frac{1}{2\times600\times\qty{24.2}{s^{-1}}} = \qty{3.4435e-5}{s}
    \quad\text{(keep the extra digits)}\\
  c &= \frac{2\times\qty{5.20e3}{m}}{\qty{3.4435e-5}{s}}
    = \frac{\qty{1.04e4}{m}}{\qty{3.4435e-5}{s}} = \qty{3.020e8}{m.s^{-1}}}}
\answer{\final{$c = \qty{3.02e8}{m.s^{-1}}$} (3 significant figures).}
\checkstep{within 1\,\% of \qty{3.00e8}{m.s^{-1}} \cmark. In \qty{34}{\micro\second} the wheel
  turns only $\frac{1}{1200}$ of a turn: a tiny turn swaps a gap for a tooth.}
\end{workedexample}

\begin{mistakes}
\begin{itemize}
  \item \xmark\ Using the one-way distance for a reflected signal. \cmark\ There and back: $d = 2D$.
  \item \xmark\ Rømer: ``the diameter of the Earth''. \cmark\ The diameter of Earth's
        \emph{orbit}.
  \item \xmark\ Fizeau: $t = \dfrac{1}{Nf}$. \cmark\ Gap to the next \emph{tooth} is half of
        that: $t = \dfrac{1}{2Nf}$.
  \item \xmark\ ``The second is defined from $c$.'' \cmark\ The second comes from caesium
        clocks; the \emph{metre} comes from $c$ and the second.
\end{itemize}
\end{mistakes}

\begin{examtip}[breakable]
\textbf{``Explain why measurements of $c$ became more accurate over time'' (3--4 marks).} Name
at least two methods with their results. For each, say \emph{what} was measured (a time, a
distance) and \emph{how well}. Give the reason for each improvement:
\begin{itemize}
  \item timing by people was replaced by machines (wheels, rotating mirrors), then by
        electronics and atomic clocks;
  \item huge, poorly known distances in space were replaced by measured paths on Earth, then by
        laser interference;
  \item many readings were repeated and averaged.
\end{itemize}
Example: Fizeau judged by eye when the light vanished behind a spinning wheel
(\qty{3.1e8}{m.s^{-1}}, about 4--5\,\% high). Michelson used a \qty{35}{km} measured path, a
fast mirror and many repeats (\qty{2.998e8}{m.s^{-1}}). End with lasers (1972), which gave $c$
to 9 significant figures, so $c$ was fixed in 1983.\\[3pt]
\textbf{``Assess the validity of Galileo's method'' (3 marks).} Judge it (invalid), give the
reason (it measured reaction time, not the travel time of light) and give numbers. Never write
just ``human error'': name the error.
\end{examtip}

\begin{quickcheck}
  \item A radar pulse reflects from an aircraft and returns after \qty{1.20e-4}{s}. How far
        away is the aircraft?
  \item Galileo's lanterns were about \qty{1.5}{km} apart. (a) Calculate the time for light to go
        there and back. (b) Explain why he could not measure it.
  \item Fizeau's wheel had 720 teeth and his mirror was \qty{8633}{m} away. Taking
        $c = \qty{3.00e8}{m.s^{-1}}$, how many turns per second make the light first disappear?
  \item State how the metre is defined today. Give two reasons why this is better than a metal
        bar.
\end{quickcheck}

%% ==================================================================
\section{Spectra and identifying elements}\label{ch1:sec:spectra}
\covers{M7.IQ1.04, M7.IQ1.05}

\begin{bigidea}
A spectrum is light spread out into its wavelengths. Hot dense objects give all colours;
low-density gases give bright lines, or dark lines when they are in front of a hotter source.
Every element has its own set of lines, like a fingerprint.
\end{bigidea}

\begin{keywords}
  \kw{Spectrum (plural spectra)}{light spread out in order of wavelength.}
  \kw{Spectroscope}{an instrument that spreads light into a spectrum so you can study it.}
  \kw{Discharge tube}{a glass tube of gas at low pressure that glows when a high voltage is
    applied.}
  \kw{Incandescent}{glowing because it is hot.}
\end{keywords}

A rainbow is a spectrum: sunlight spread into its colours. A \term{spectroscope} does this with
a prism or a \term{diffraction grating}, a plate with thousands of tiny, equal slits
(\cref{ch2:chap}). A grating bends red light the most; a prism bends violet the most.
\Cref{ch1:fig:spectroscope} (Key experiments) shows one. Using spectra to learn about a light
source is called \term{spectroscopy}.

\subsection{Three kinds of spectra}

\begin{enumerate}
  \item A \term{continuous spectrum} has all wavelengths, with no gaps. It comes from a hot,
        dense object: an \term{incandescent} filament, lava, the inside of a star.
  \item An \term{emission spectrum} has bright lines on a dark background. It comes from a
        \term{low-density} gas: a gas whose atoms are far apart. The atoms must have extra
        energy. Atoms with extra energy are \term{excited}. A flame excites atoms by heating
        them. A \term{discharge tube} or a neon sign uses a high voltage. The voltage speeds up
        free electrons, which crash into gas atoms and give them extra energy.
  \item An \term{absorption spectrum} is a continuous spectrum with dark lines. It forms when
        light from a hot dense source passes through a cooler, low-density gas. The dark lines
        are at \emph{exactly the same wavelengths} as the bright lines of that gas.
\end{enumerate}
\Cref{ch1:fig:spectra} shows the three kinds. \Cref{ch1:fig:kirchhoff} shows how each is made
(Kirchhoff's three rules, 1859).

\begin{figure}
  \centering
  \includegraphics{ch1_spectra.pdf}
  \caption{The three kinds of spectra.}\label{ch1:fig:spectra}
\end{figure}

\begin{figure}
  \centering
  \begin{tikzpicture}[x=1cm, y=1cm]
    \draw[line width=0.8pt, fill=hscray!25] (0,0) circle (0.5);
    \draw[line width=0.9pt, hscacc] (-0.28,-0.05) -- (-0.19,0.14) -- (-0.09,-0.14) -- (0,0.14)
      -- (0.09,-0.14) -- (0.19,0.14) -- (0.28,-0.05);
    \node[diagramlabel, anchor=east, align=right] at (-0.6,0) {hot, dense\\source};
    \draw[line width=0.7pt, hscmuted, fill=hscmuted!12] (3.5,0) ellipse (1.2 and 0.86);
    \node[diagramlabel, align=center] at (3.5,0) {cool,\\low-density\\gas};
    \draw[ray] (0.55,0) -- (2.28,0);
    \draw[ray] (4.72,0) -- (6.3,0);
    \draw[ray] (0,-0.55) -- (0,-1.2);
    \draw[ray] (3.5,-0.88) -- (3.5,-1.25);
    \shade[shading=ch1rainbow] (-1.1,-1.7) rectangle (1.1,-1.3);
    \node[diagramlabel, anchor=north] at (0,-1.75) {continuous};
    \fill[black] (2.4,-1.7) rectangle (4.6,-1.3);
    \foreach \x/\col in {0.073/violet, 0.249/blue, 0.631/cyan, 1.877/red}
      \fill[\col] ($(2.4,-1.7)+(\x-0.025,0)$) rectangle ($(2.4,-1.3)+(\x+0.025,0)$);
    \node[diagramlabel, anchor=north, align=center] at (3.5,-1.75) {emission\\(seen from the side)};
    \shade[shading=ch1rainbow] (6.4,-0.2) rectangle (8.6,0.2);
    \foreach \x in {0.073, 0.249, 0.631, 1.877}
      \fill[black] ($(6.4,-0.2)+(\x-0.025,0)$) rectangle ($(6.4,0.2)+(\x+0.025,0)$);
    \node[diagramlabel, anchor=north] at (7.5,-0.25) {absorption};
  \end{tikzpicture}
  \caption{Kirchhoff's rules. The gas absorbs some wavelengths from the beam and gives them out
    again in all directions.}\label{ch1:fig:kirchhoff}
\end{figure}

\subsection{Why lines? Energy levels}

\begin{sidebyside}[0.40]
  \centering
  \begin{tikzpicture}[x=1cm, y=1cm]
    \draw[->, axisline] (-0.35,0) -- (-0.35,2.7) node[above, font=\footnotesize] {energy};
    \foreach \y/\n in {0/1, 1.45/2, 2.1/3, 2.45/4}
      \draw[line width=0.9pt, hscink] (0,\y) -- (3.0,\y) node[right, font=\footnotesize] {$n=\n$};
    \draw[->, line width=1.1pt, hscforce] (0.95,2.1) -- (0.95,1.48);
    \draw[->, decorate, decoration={snake, amplitude=1.3pt, segment length=5pt, post length=2mm},
      line width=0.8pt, hscray] (0.85,1.78) -- (0.05,1.78);
    \node[font=\footnotesize, align=center, anchor=north] at (0.5,1.38) {emission:\\photon out};
    \draw[->, line width=1.1pt, hscvel] (2.05,1.45) -- (2.05,2.07);
    \draw[->, decorate, decoration={snake, amplitude=1.3pt, segment length=5pt, post length=2mm},
      line width=0.8pt, hscray] (2.95,1.78) -- (2.15,1.78);
    \node[font=\footnotesize, align=center, anchor=north] at (2.5,1.38) {absorption:\\photon in};
    \node[font=\footnotesize, anchor=north] at (1.5,-0.05) {lowest level};
  \end{tikzpicture}
  \captionof{figure}{Electrons jump between energy levels (straight arrows). The wavy arrows are
    photons.}\label{ch1:fig:levels}
\nextside
  \begin{itemize}
    \item \Cref{ch1:fig:levels} shows that electrons in an atom can only have certain energies,
          called \term{energy levels}.
    \item An electron that drops to a lower level gives out a \term{photon}. A photon is a
          packet of light energy. Its energy equals the gap between the two levels. One gap
          gives one frequency ($E = hf$, where $h$ is Planck's constant; \cref{ch3:chap}), so
          one bright line.
    \item An atom absorbs only photons whose energy matches a gap. Those wavelengths go missing
          from the light that passes through, so we see dark lines.
  \end{itemize}
\end{sidebyside}

Each element has its own energy levels, so its own set of lines. Its emission and absorption
lines come from the same gaps, so they are at the same wavelengths.

\subsection{Identifying elements}

To identify the elements in a sample: measure the wavelengths of its lines, and compare them with
laboratory reference lines (\cref{ch1:tab:lines}). An element is present only if \emph{all} of
its strong lines are present. One matching line is not enough. If a line is given as a
frequency, change it to a wavelength with $c = f\lambda$. (There is no new formula here.)

\begin{table}
  \caption{The strongest visible lines of four elements.}\label{ch1:tab:lines}
  \centering
  \begin{tabular}{@{}L{30mm}L{90mm}@{}}
    \toprule
    \colhead{Element} & \colhead{Bright lines (\si{nm})}\\
    \midrule
    Hydrogen (H) & 410, \ 434, \ 486, \ 656\\
    Helium (He) & 447, \ 471, \ 492, \ 502, \ 588, \ 668\\
    Sodium (Na) & 589 (two lines very close together)\\
    Mercury (Hg) & 405, \ 436, \ 546, \ 577, \ 579\\
    \bottomrule
  \end{tabular}
\end{table}

Chemists use this to test mineral samples and check for pollution. Helium was discovered in
1868 in light from the Sun. Scientists saw a yellow line that matched no known element. Helium
was found on Earth 27 years later.

\begin{workedexample}{Which kind of spectrum?}\label{ch1:we:kirchhoff}
Predict the kind of spectrum, with a reason: (a) the hot filament of a car headlight; (b) a
low-pressure sodium street lamp; (c) light from a filament lamp after it passes through cool
sodium gas.
\given{three light sources}
\find{the kind of spectrum for each}
\formula{Kirchhoff's rules: hot and dense $\to$ continuous; hot or excited low-density gas
  $\to$ emission; cool low-density gas in front of a hot source $\to$ absorption}
\substitute{(a) a hot solid, so continuous. (b) sodium gas at low pressure, excited by a high
  voltage, so emission.
  (c) continuous light through cooler, low-density sodium gas, so absorption.}
\answer{(a) continuous; (b) emission, a bright yellow line at \qty{589}{nm}; (c) absorption: all
  colours with a dark line at \qty{589}{nm}.}
\checkstep{the dark line in (c) is where the bright line in (b) is: same gas, same energy gaps
  \cmark.}
\end{workedexample}

\begin{workedexample}{Identifying the gases in a tube}\label{ch1:we:identify}
A discharge tube holds a mixture of gases. Its lines are at 405, 410, 436, 486, 546, 577, 579
and \qty{656}{nm}, and one more line has a frequency of \qty{6.91e14}{Hz}. Identify the gases.
\given{eight wavelengths; one line with $f = \qty{6.91e14}{Hz}$; $c = \qty{3.00e8}{m.s^{-1}}$}
\find{the elements in the tube}
\formula{$\lambda = \dfrac{c}{f}$, then: all of an element's lines must be present}
\substitute{$\lambda = \dfrac{\qty{3.00e8}{m.s^{-1}}}{\qty{6.91e14}{Hz}} = \qty{4.34e-7}{m}
  = \qty{434}{nm}$.\\[6pt]
  H: 410, 434, 486, 656 \cmark\ all there. Hg: 405, 436, 546, 577, 579 \cmark\ all there.\\
  He: no 588 \xmark. Na: no 589 \xmark.}
\answer{the tube contains \final{hydrogen and mercury}.}
\checkstep{every measured line belongs to hydrogen or mercury. None is left over \cmark.}
\end{workedexample}

\begin{mistakes}
\begin{itemize}
  \item \xmark\ ``A hot filament gives an emission spectrum.'' \cmark\ A filament is a hot, dense
        solid: continuous.
  \item \xmark\ ``Absorption lines are at different wavelengths from emission lines.'' \cmark\
        For the same element they are at the same wavelengths.
  \item \xmark\ ``One matching line proves the element is there.'' \cmark\ All its strong lines
        must be there.
\end{itemize}
\end{mistakes}

\begin{examtip}
\textbf{``Contrast emission and absorption spectra'' (3 marks).} ``Contrast'' means give
differences. (1) How each is made: an excited low-density gas emits; a cool gas absorbs from
continuous light passing through it. (2) How each looks: bright lines on black; dark lines on a
continuous spectrum. (3) Where you look from: emission is seen from any direction; absorption
only when you look through the gas at the hot source. You may add the link (same element, same
wavelengths, same energy gaps), but the marks are for the differences.\\[3pt]
\textbf{``How does spectroscopy identify elements?'' (2 marks).} Each element has a unique set of
lines; compare with reference spectra; all lines must match.
\end{examtip}

\begin{quickcheck}
  \item Name the three kinds of spectra. Give one source of each.
  \item Explain why a gas has its dark absorption lines at the same wavelengths as its bright
        emission lines.
  \item A gas gives lines at 405, 436, 447, 471, 492, 502, 546, 577, 579, 588 and \qty{668}{nm}.
        Use \cref{ch1:tab:lines} to identify the elements.
  \item A student sees one line at \qty{486}{nm} and concludes that hydrogen is present.
        Evaluate her conclusion.
\end{quickcheck}

%% ==================================================================
\section{What starlight tells us}\label{ch1:sec:stars}
\covers{M7.IQ1.06}

\begin{bigidea}
A star's spectrum is a continuous spectrum crossed by dark absorption lines. The peak gives the
surface temperature, and the lines present give the chemical composition. Shifted lines show
motion towards or away from us; wider lines show fast rotation or high density.
\end{bigidea}

\begin{keywords}
  \kw{Black body}{an ideal object that absorbs all radiation and is a perfect emitter. Its
    spectrum depends only on its temperature.}
  \kw{Doppler effect}{the change in wavelength when a source and an observer move towards or
    away from each other.}
  \kw{Redshift, blueshift}{lines moved to longer (red) or shorter (blue) wavelengths.}
  \kw{Translational velocity}{how fast the whole star moves through space.}
  \kw{Rotational velocity}{how fast the star spins (the speed of its surface as it turns).}
  \kw{Line broadening}{spectral lines becoming wider.}
\end{keywords}

The inside of a star is hot and dense, so it gives a continuous spectrum. This light passes
through the star's cooler, lower-density outer layers, called its \term{atmosphere}. Atoms there
absorb the wavelengths that match their own energy gaps (\cref{ch1:sec:spectra}). So a star
gives an absorption spectrum. The dark lines in sunlight are called \term{Fraunhofer lines}.

\textbf{1. Chemical composition.} Match the dark lines with the laboratory lines of each element
(\cref{ch1:sec:spectra}). The Sun's dark lines show hydrogen, sodium, magnesium, calcium, iron and
many more elements.

\textbf{2. Surface temperature.} The continuous part is close to a \term{black body} curve
(\cref{ch1:fig:blackbody}; more in \cref{ch3:sec:blackbody}). A hotter star peaks at a shorter wavelength (Wien's law). It gives
out more light at every wavelength from each square metre of its surface, and it looks bluer.
You can see the same thing with a metal rod in a fire. As it gets hotter, it glows dull red,
then orange, then yellow-white.

Which lines are strong also depends on temperature. Stars near \qty{10000}{K} have the strongest
hydrogen lines. Cool stars (below about \qty{3700}{K}) show wide dark bands made by molecules.
Astronomers sort stars into classes O, B, A, F, G, K, M, from hottest to coolest (Module 8).

\begin{figure}
  \centering
  \includegraphics{ch1_blackbody.pdf}
  \caption{Black-body spectra at three temperatures.}\label{ch1:fig:blackbody}
\end{figure}

\begin{refresher}{Temperature in kelvin (Year 11)}
Physics formulas use absolute temperature $T$ in kelvin (\si{K}).
$T(\unit{K}) = T(\unit{\degreeCelsius}) + 273$. Example: \qty{27}{\degreeCelsius} $=$
\qty{300}{K}. Kelvin has no degree sign.
\end{refresher}

\begin{formulacard}{Wien's displacement law}
  \cardformula{\lambda_\mathrm{max} = \frac{b}{T}}
  \begin{symbols}
    \sym{$\lambda_\mathrm{max}$}{peak wavelength of the continuous spectrum}{\si{m}}
    \sym{$b$}{Wien's displacement constant, \num{2.898e-3} (data sheet)}{\si{m.K}}
    \sym{$T$}{surface temperature}{\si{K}}
  \end{symbols}
  \inwords{hotter objects give out most of their light at shorter wavelengths.}
  \whentouse{the peak of the continuous (black-body) spectrum of a star or a hot object.}
  \whennottouse{on the wavelength of one emission or absorption line. A line is not the peak.}
  \rearranged{$T = \dfrac{b}{\lambda_\mathrm{max}}$ \qquad $b = \lambda_\mathrm{max}T$}
  \tensecond{the Sun's surface is at about \qty{5800}{K}, so
    $\lambda_\mathrm{max} = \dfrac{2.898\times10^{-3}}{5800} = \qty{5.00e-7}{m} = \qty{500}{nm}$.}
  \trap{$T$ must be in kelvin: add 273 to a temperature in \si{\degreeCelsius}. Change \si{nm}
    to \si{m}.}
\end{formulacard}

\textbf{3. Translational velocity.} \term{Translational velocity} is how fast the whole star
moves through space. A siren sounds higher as it comes towards you and lower as it goes away:
this is the \term{Doppler effect}. Light does the same (\cref{ch1:fig:starlines}).

A star moving away gives a \term{redshift}: every line moves to a longer wavelength. A star
moving towards us gives a \term{blueshift}. The whole pattern moves together, so the elements
can still be recognised. A faster star gives a bigger shift.

Only motion along our \term{line of sight} (the straight line from us to the star) shows up:
motion straight towards or away from us. Each line moves by the fraction $v/c$ of its
wavelength, where $v$ is the speed along the line of sight: $\Delta\lambda = \dfrac{\lambda v}{c}$
\hscbadgenotonsheet. Where it comes from: in one period, $T = 1/f = \lambda/c$ (here $T$ is the
period), a star moving away travels $vT$. So its crests arrive $\lambda + vT$ apart, and $\Delta\lambda = vT =
\lambda v/c$. This works only when $v$ is much smaller than $c$.

The shifts are tiny. For $v = \qty{30}{km.s^{-1}} = \qty{3.0e4}{m.s^{-1}}$, the \qty{656}{nm}
line moves by only $\Delta\lambda = (656\times10^{-9})\times(3.0\times10^{4}) \div
(3.00\times10^{8}) = \qty{6.6e-11}{m}$, about \qty{0.07}{nm}. So spectroscopes must be very
precise. HSC questions usually ask about shifts in words.

\begin{figure}
  \centering
  \includegraphics{ch1_starlines.pdf}
  \caption{Hydrogen absorption lines in four cases.}\label{ch1:fig:starlines}
\end{figure}

\begin{sidebyside}[0.56]
  \centering
  \begin{tikzpicture}[x=1cm, y=1cm]
    \draw[line width=0.8pt, fill=hscray!20] (0,0) circle (0.85);
    \draw[->, line width=0.8pt, hscmuted] (60:1.08) arc (60:120:1.08);
    \node[font=\footnotesize, text=hscmuted, anchor=south] at (0,1.08) {spin (top view)};
    \draw[->, line width=1.3pt, hscvel] (-0.85,0.25) -- (-0.85,-0.7);
    \node[font=\footnotesize, align=right, anchor=east] at (-0.95,-0.25) {towards us:\\blueshift};
    \draw[->, line width=1.3pt, hscforce] (0.85,-0.25) -- (0.85,0.7);
    \node[font=\footnotesize, align=left, anchor=west] at (0.95,0.25) {away:\\redshift};
    \draw[->, line width=0.6pt, hscmuted] (0,-0.95) -- (0,-1.35) node[below, font=\footnotesize] {to Earth};
    \begin{scope}[xshift=2.55cm, yshift=-0.9cm]
      \draw[axisline] (0,0) -- (2.1,0) node[below, font=\footnotesize] {$\lambda$};
      \draw[axisline] (0,0) -- (0,1.85) node[left, font=\footnotesize] {light};
      \draw[hscseriesA, line width=1pt, domain=0.05:2.0, samples=120, smooth]
        plot (\x, {1.55 - 1.15*exp(-((\x-1.05)/0.1)^2)});
      \draw[hscseriesB, line width=1pt, dashed, domain=0.05:2.0, samples=120, smooth]
        plot (\x, {1.55 - 0.29*exp(-((\x-1.05)/0.4)^2)});
      \node[font=\footnotesize, text=hscseriesA, anchor=west] at (1.17,0.42) {slow};
      \node[font=\footnotesize, text=hscseriesB, anchor=south] at (1.05,1.58) {fast};
    \end{scope}
  \end{tikzpicture}
  \captionof{figure}{A rotating star, and one line for a slow and a fast rotator.}
  \label{ch1:fig:rotation}
\nextside
  \textbf{4. Rotational velocity.} \term{Rotational velocity} is how fast the star spins (the
  speed of its surface as it turns). One edge of a rotating star turns towards us (blueshift).
  The other edge turns away (redshift). We see the whole star at once, so each line is spread
  out: it is \term{broadened}. Its centre does not move (\cref{ch1:fig:rotation}). Faster
  rotation gives wider lines.
\end{sidebyside}

\textbf{5. Density.} In a dense atmosphere, atoms collide often and sit close to other charged
particles. This disturbs their energy levels a little, so each atom absorbs a slightly different
wavelength. The lines get wider (\term{pressure broadening}).

\term{Giant stars} are huge, swollen stars near the end of their lives. Their outer layers are
spread out, so the density is low and the lines are narrow. \term{Main-sequence stars}, such as
the Sun, are ordinary stars in the long, middle part of their lives. Their outer layers are
denser, so their lines are broader (Module 8). Rotation and density both widen lines; a question
will say which one changes. \Cref{ch1:tab:readstar} sums up how to read a star's spectrum.

\begin{table}
  \caption{Reading a star's spectrum.}\label{ch1:tab:readstar}
  \centering
  \begin{tabular}{@{}L{72mm}L{68mm}@{}}
    \toprule
    \colhead{What you see} & \colhead{What it tells you}\\
    \midrule
    which lines are present & chemical composition\\
    peak of the continuous curve; which lines are strong & surface temperature\\
    whole pattern shifted to red or to blue & moving away or towards (translational velocity)\\
    lines wider, centres not moved & faster rotation, or higher density\\
    \bottomrule
  \end{tabular}
\end{table}

\begin{workedexample}{The temperature of a star}\label{ch1:we:wien}
A star gives out the most light at a wavelength of \qty{414}{nm}. Calculate its surface
temperature.
\given{$\lambda_\mathrm{max} = \qty{414}{nm} = \qty{4.14e-7}{m}$, \ $b = \qty{2.898e-3}{m.K}$}
\find{the surface temperature $T$}
\formula{$\lambda_\mathrm{max} = \dfrac{b}{T}$. Multiply both sides by $T$, then divide both
  sides by $\lambda_\mathrm{max}$: $T = \dfrac{b}{\lambda_\mathrm{max}}$}
\substitute{$T = \dfrac{\qty{2.898e-3}{m.K}}{\qty{4.14e-7}{m}} = \qty{7.00e3}{K}$}
\answer{\final{$T = \qty{7.00e3}{K}$} (3 significant figures).}
\checkstep{the peak is bluer than the Sun's (\qty{500}{nm}), so the star should be hotter than
  \qty{5800}{K}. It is \cmark.}
\end{workedexample}

\begin{workedexample}{Two stars compared}\label{ch1:we:twostars}
The table gives four hydrogen absorption lines (\si{nm}). (a) What do the spectra show about the
motion of each star towards or away from Earth? (b) Give two possible reasons why the lines of
star Q are wider.
\begin{center}\small
  \begin{tabular}{@{}lccccl@{}}
    \toprule
    Laboratory & 410.2 & 434.0 & 486.1 & 656.3 & \\
    Star P & 410.5 & 434.3 & 486.5 & 656.8 & narrow lines\\
    Star Q & 410.2 & 434.0 & 486.1 & 656.3 & wide lines\\
    \bottomrule
  \end{tabular}
\end{center}
\given{laboratory and star wavelengths; line widths}
\find{(a) the motion of P and Q; (b) reasons for Q's wide lines}
\formula{rules: whole pattern shifted $\to$ moving (longer = away); wider lines $\to$ rotation or
  density}
\substitute{(a) P: every line is longer, by 0.3 to \qty{0.5}{nm}: a redshift. Q: no shift.
  (b) Q's lines are wide but not moved.}
\answer{(a) P is moving \final{away from Earth}; Q is not moving towards or away (it may move
  across the sky). (b) Q rotates faster, or has a denser atmosphere.}
\checkstep{each of P's lines moves by about the same fraction of its wavelength (0.069 to
  0.082\,\%; the spread comes from rounding to \qty{0.1}{nm}). The longest line (\qty{656}{nm})
  moves the most, as a Doppler shift should \cmark.}
\end{workedexample}

\begin{mistakes}
\begin{itemize}
  \item \xmark\ Using $\lambda_\mathrm{max} = \dfrac{b}{T}$ on a dark line. \cmark\ Use the peak
        of the continuous curve, with $T$ in kelvin.
  \item \xmark\ ``Redshift means the star looks red.'' \cmark\ It means the lines move to longer
        wavelengths. A blue star can be redshifted.
  \item \xmark\ Mixing up shift and broadening. \cmark\ Motion moves the lines. Rotation and
        density widen them.
  \item \xmark\ Using the sheet's Doppler formula for light. \cmark\ It is for sound. Describe
        light shifts in words.
\end{itemize}
\end{mistakes}

\begin{examtip}
\textbf{``Explain how a star's spectrum gives information about \ldots'' (5--6 marks)} or
\textbf{``Evaluate how useful spectroscopy is to astronomers''.} Start: a star's spectrum is
continuous with dark lines (hot dense inside, cooler outer gas). Then one short paragraph per
property, with a heading: what you see $\to$ why $\to$ what you conclude. For ``evaluate'', also
give a limit (only motion along the line of sight shows up; rotation and density both widen
lines). End with a judgement: spectroscopy is very useful because all this comes from starlight
alone.\\[3pt]
\textbf{Multiple choice:} same line positions but narrower lines means lower density (or slower
rotation).
\end{examtip}

\begin{quickcheck}
  \item A red star's continuous spectrum peaks at \qty{828}{nm}. Calculate its surface
        temperature.
  \item A lamp filament is at \qty{2527}{\degreeCelsius}. (a) Calculate its peak wavelength.
        (b) Explain why filament lamps waste a lot of energy.
  \item Explain how a star's spectrum shows that it rotates. Why does the centre of each line
        not move?
  \item Two stars have their lines at the same wavelengths, and their surfaces move at the same
        speed as they rotate. Star A has narrow lines; star B has broad lines. Which has the
        denser atmosphere? Explain.
  %% Added by the Module 7 coverage audit (M7.IQ1.06, translational velocity).
  \item In a star's spectrum a hydrogen line is at \qty{656.17}{nm}. In the laboratory the same
        line is at \qty{656.30}{nm}. All the star's other lines move by the same fraction.
        (a) Is the star moving towards Earth or away from it? Explain. (b) Use
        $\Delta\lambda = \dfrac{\lambda v}{c}$ to estimate its speed along our line of sight.
        (c) Explain why the star's real speed through space may be bigger.\label{ch1:qc:doppler}
\end{quickcheck}

%% ==================================================================
% Layout (book QA): keep the heading and its two lines with the first experiment box.
\Needspace{16\baselineskip}
\section{Key experiments}\label{ch1:sec:experiments}
\covers{M7.IQ1.03, M7.IQ1.04}

The Working scientifically words box in \cref{ch1:sec:speed} explains validity, reliability
and accuracy.

\begin{experiment}{Galileo's lanterns (about 1638)}
\aim{to find out whether light takes time to travel.}
\apparatus{two covered lanterns, two people on hills about \qty{1.5}{km} apart, a simple timer.
  \begin{figure}
    \centering
    \begin{tikzpicture}[x=1cm, y=1cm]
      \fill[hscmuted!25] (-0.4,0) .. controls (0.3,0.75) and (1.1,0.75) .. (1.8,0) -- cycle;
      \fill[hscmuted!25] (5.2,0) .. controls (5.9,0.75) and (6.7,0.75) .. (7.4,0) -- cycle;
      \fill[hscray] (0.7,0.66) circle (0.11);
      \fill[hscray] (6.3,0.66) circle (0.11);
      \draw[ray] (0.95,0.8) -- (6.05,0.8);
      \draw[ray] (6.05,0.56) -- (0.95,0.56);
      \node[font=\footnotesize, anchor=south] at (3.5,0.83) {light out};
      \node[font=\footnotesize, anchor=north] at (3.5,0.53) {light back};
      \node[font=\footnotesize, anchor=north] at (0.7,-0.02) {Galileo (timer)};
      \node[font=\footnotesize, anchor=north] at (6.3,-0.02) {assistant};
    \end{tikzpicture}
    \caption{Galileo's lantern experiment.}\label{ch1:fig:galileo}
  \end{figure}}
\method{Galileo uncovers his lantern and starts his timer. His assistant uncovers his own
  lantern as soon as he sees Galileo's light. Galileo stops the timer when he sees the
  assistant's light (\cref{ch1:fig:galileo}). Repeat with the lanterns close together and
  compare the times.}
\results{the times were the same for near and far lanterns. No delay was measured.}
\conclusion{light is either instant or far too fast to measure this way.}
\significance{an early test by experiment. A method must suit the size of what it measures.}
\validity{\textbf{Validity:} invalid. Light goes \qty{3}{km} there and back in
  \qty{1.0e-5}{s}, but reaction time (about \qty{0.2}{s}) is \num{20000} times longer. It
  measured reaction time. \textbf{Reliability:} repeats give similar times (all reaction time):
  reliable, but still invalid. \textbf{Accuracy:} no value for $c$ at all.}
\end{experiment}

\begin{experiment}{Rømer and the moons of Jupiter (1676)}
\aim{to explain why the eclipses of Jupiter's moon Io do not keep to a regular timetable.}
\apparatus{telescope, pendulum clock, tables of predicted eclipse times.
  \begin{figure}
    \centering
    \begin{tikzpicture}[x=1cm, y=1cm]
      \draw[guide] (0,0) circle (1.4);
      \fill[hscray] (0,0) circle (0.22);
      \node[font=\footnotesize, anchor=north] at (0,-0.24) {Sun};
      \coordinate (Enear) at (-14:1.4);
      \coordinate (Efar) at (166:1.4);
      \coordinate (J) at (6.0,0.12);
      \draw[ray] (J) -- (Enear);
      \draw[ray] (J) -- (Efar);
      \fill[hscvel] (Enear) circle (0.1);
      \fill[hscvel] (Efar) circle (0.1);
      \node[font=\footnotesize, align=left, anchor=north west] at ($(Enear)+(0.12,-0.08)$)
        {Earth near:\\eclipses on time};
      \node[font=\footnotesize, align=right, anchor=south east] at ($(Efar)+(-0.1,0.1)$)
        {Earth far:\\eclipses late};
      \fill[hscmuted!30] (6.0,0.4) rectangle (7.5,-0.16);
      \fill[hscacc!70!black] (J) circle (0.28);
      \draw[guide] (J) circle (0.7);
      \fill[hscink] ($(J)+(0.7,0)$) circle (0.07);
      \node[font=\footnotesize, anchor=south] at ($(J)+(0,0.72)$) {Jupiter};
      \node[font=\footnotesize, anchor=west] at ($(J)+(1.55,0)$) {Io in shadow};
      \draw[line width=0.6pt, hscink] (0,0) -- (225:1.4)
        node[pos=0.55, above left=-1pt, font=\footnotesize] {$r$};
      \draw[guide] (Enear) -- (1.358,-1.8);
      \draw[guide] (Efar) -- (-1.358,-1.8);
      \draw[<->, line width=0.6pt, hscink] (-1.358,-1.7) -- (1.358,-1.7)
        node[midway, below, font=\footnotesize] {extra distance $\approx 2r$};
      \node[font=\footnotesize, text=hscmuted] at (6.3,-1.6) {not to scale};
    \end{tikzpicture}
    \caption{Rømer's method. From the far side of its orbit, Earth receives the light of each
      eclipse later.}\label{ch1:fig:romer}
  \end{figure}}
\method{Io orbits Jupiter every 1.77 days and passes into Jupiter's shadow each orbit. Time
  eclipses when Earth is near Jupiter and use them to predict later ones. Compare predicted and
  observed times as Earth moves to the far side of its orbit (\cref{ch1:fig:romer}).}
\results{the eclipses came later and later, up to about \qty{22}{min} late (modern value
  \qty{16.6}{min}), then earlier again as Earth came back.}
\conclusion{light has a finite speed. The extra delay is the time to cross Earth's orbit.
  Huygens used it to get $c \approx \qty{2.2e8}{m.s^{-1}}$.}
\significance{the first evidence that light is not instant, and the first estimate of its speed.}
\validity{\textbf{Validity:} valid in principle. Io's orbit is a very regular natural clock.
  The only thing that changes is the Earth--Jupiter distance. So the extra delay measures the
  extra distance that light travels. The extra distance is huge, so the delay is long enough to
  measure. \textbf{Reliability:} timing many eclipses over many months reduces random error.
  \textbf{Accuracy:} Huygens' value, \qty{2.2e8}{m.s^{-1}}, is about 27\,\% too low
  ($(3.00 - 2.2) \div 3.00 \times 100\,\% \approx 27\,\%$). Io fades into the shadow slowly,
  clocks were poor, and the size of Earth's orbit was poorly known. Using $2r$ for the extra
  distance is also an approximation. Earth's orbit is not a perfect circle, and Jupiter moves a
  little in six months. If the distance used is smaller than the real extra distance, the value
  of $c$ comes out too small.}
\end{experiment}

\begin{experiment}{Fizeau's toothed wheel (1849)}
\aim{to measure the speed of light directly, on Earth.}
\apparatus{lamp; a half-silvered mirror (a glass plate that reflects about half the light and
  lets the rest through); a toothed wheel (720 teeth) turned at a known rate; a mirror
  \qty{8.63}{km} away (in Paris); an eyepiece (the lens the observer looks through).
  \begin{figure}
    \centering
    \begin{tikzpicture}[x=1cm, y=1cm]
      \draw[line width=0.8pt] (0,0) circle (0.2);
      \fill[hscink] (0.09,0) circle (0.07);
      \node[font=\footnotesize, anchor=north] at (0,-0.25) {observer};
      \draw[line width=1.6pt, hscfield!80] (1.1,-0.28) -- (1.66,0.28);
      \node[font=\footnotesize, align=center, anchor=south] at (1.38,0.35) {half-silvered\\mirror};
      \draw[line width=0.8pt, fill=hscray!35] (1.38,-1.3) circle (0.2);
      \node[font=\footnotesize, anchor=east] at (1.12,-1.3) {lamp};
      \draw[ray] (1.38,-1.07) -- (1.38,-0.07);
      \foreach \y in {-0.95,-0.55,0.25,0.65}
        \fill[hscink] (2.9,\y) rectangle (3.08,\y+0.3);
      \draw[line width=0.6pt, hscmuted] (2.99,-1.05) -- (2.99,1.05);
      \node[font=\footnotesize, align=left, anchor=west] at (3.15,-0.85) {wheel\\(edge-on)};
      \draw[ray] (1.44,0.06) -- (4.5,0.06);
      \draw[ray] (5.7,0.06) -- (8.1,0.06);
      \draw[ray] (8.1,-0.08) -- (5.7,-0.08);
      \draw[ray] (4.5,-0.08) -- (0.28,-0.08);
      \draw[line width=0.6pt] (4.85,-0.28) -- (5.05,0.28);
      \draw[line width=0.6pt] (5.15,-0.28) -- (5.35,0.28);
      \draw[line width=2pt, hscink] (8.2,-0.55) -- (8.2,0.55);
      \fill[ground] (8.22,-0.55) rectangle (8.42,0.55);
      \node[font=\footnotesize, anchor=north] at (8.25,-0.6) {mirror};
      \draw[<->, line width=0.6pt, hscmuted] (2.99,-1.6) -- (8.2,-1.6)
        node[midway, fill=white, font=\footnotesize, text=hscmuted] {$D = \qty{8.63}{km}$};
            \begin{scope}[shift={(6.1,1.15)}]
        \foreach \i in {0,...,15}
          \fill[hscink] ({22.5*\i-5.6}:0.45) -- ({22.5*\i-5.6}:0.62)
            arc ({22.5*\i-5.6}:{22.5*\i+5.6}:0.62) -- ({22.5*\i+5.6}:0.45)
            arc ({22.5*\i+5.6}:{22.5*\i-5.6}:0.45) -- cycle;
        \draw[line width=0.8pt, fill=hscmuted!20] (0,0) circle (0.45);
        \fill[hscink] (0,0) circle (0.05);
        \node[font=\footnotesize, align=left, anchor=west] at (0.72,0) {wheel face-on:\\$N$ teeth, $N$ gaps};
      \end{scope}
    \end{tikzpicture}
    \caption{Fizeau's apparatus. Light leaves through a gap. At the right speed, it returns to
      find a tooth.}\label{ch1:fig:fizeau}
  \end{figure}}
\method{light reflects off the half-silvered mirror, passes through a gap in the wheel, goes to
  the far mirror and back (\cref{ch1:fig:fizeau}). Spin the wheel faster until the returning
  light first hits the next tooth and the observer sees darkness. Record the rate.}
\results{darkness at about 12.6 turns per second, so $t = \dfrac{1}{2\times720\times12.6} =
  \qty{5.51e-5}{s}$ and $c = \dfrac{2\times\qty{8633}{m}}{\qty{5.51e-5}{s}} =
  \qty{3.13e8}{m.s^{-1}}$.}
\conclusion{$c \approx \qty{3.1e8}{m.s^{-1}}$, within about 5\,\% of today's value.}
\significance{the first measurement of $c$ on Earth. Maxwell compared his prediction with values
  like this.}
\validity{\textbf{Validity:} a direct method, with a measured distance and a measured time. A
  distance on Earth can be surveyed accurately. \textbf{Reliability:} repeat many times. Also
  spin the wheel faster. The light disappears again at 3 times and 5 times the first speed (the
  light then meets the 2nd or 3rd tooth). This checks the first result. \textbf{Accuracy:} about 4--5\,\% high. As the wheel speeds up, the
  returning light gets dimmer gradually. So it is hard to judge by eye the exact speed at which
  it disappears. The wheel's speed must also stay steady.}
\end{experiment}

\begin{experiment}{Measuring $c$ with a microwave oven (a modern school method)}
\aim{to find $c$ from $c = f\lambda$ (an indirect method).}
\apparatus{microwave oven with the turntable removed (so the chocolate stays still), a flat
  bar of chocolate, ruler. The oven's label gives $f = \qty{2.45}{GHz}$.
  \begin{figure}
    \centering
    \begin{tikzpicture}[x=1cm, y=1cm]
      \draw[line width=0.7pt, fill=hscacc!35!brown!60] (0,0) rectangle (6.2,0.7);
      \foreach \x in {0.85,2.35,3.85,5.35}
        \fill[hscacc!20!brown!35, draw=hscacc!50!black, line width=0.4pt] (\x,0.35) ellipse (0.3 and 0.25);
      \draw[<->, line width=0.6pt, hscink] (2.35,0.95) -- (3.85,0.95)
        node[midway, above, font=\footnotesize] {$\lambda/2$};
      \draw[guide] (2.35,0.35) -- (2.35,1.05);
      \draw[guide] (3.85,0.35) -- (3.85,1.05);
      \node[font=\footnotesize, anchor=west] at (6.35,0.35) {melted spots};
    \end{tikzpicture}
    \caption{Melted spots are half a wavelength apart.}\label{ch1:fig:microwave}
  \end{figure}}
\method{heat the chocolate for about \qty{20}{s}, until a few spots just melt. Measure the
  distance between the centres of neighbouring melted spots (\cref{ch1:fig:microwave}). Repeat
  with fresh chocolate.}
\results{the microwaves reflect off the metal walls of the oven. The reflected waves overlap the
  incoming waves and make a \term{standing wave}: a wave pattern that stays in one place
  instead of travelling. At some fixed spots, called \term{antinodes}, the wave is always
  strong, so the chocolate melts there. Neighbouring antinodes are half a wavelength apart. When
  one is at a crest, the next one is at a trough. Both heat the chocolate equally, because
  heating does not depend on the direction of the field. So
  $\lambda = 2\times\text{(spacing)}$. Example: spacings 6.0, 6.3 and \qty{6.2}{cm}, mean
  \qty{6.17}{cm}. So
  $\lambda = 2\times\qty{6.167}{cm} = \qty{0.1233}{m}$ and
  $c = f\lambda = \qty{2.45e9}{Hz}\times\qty{0.1233}{m} = \qty{3.02e8}{m.s^{-1}}$.}
\conclusion{$c \approx \qty{3.0e8}{m.s^{-1}}$, within about 1\,\% of the accepted value.}
\significance{an indirect method: $c$ from $f$ and $\lambda$, not from timing. In 1972 the same
  idea with a laser gave $c$ to 9 significant figures. This led to $c$ being fixed in 1983.}
\validity{\textbf{Validity:} remove the turntable so the chocolate stays still in the
  standing-wave pattern. The method assumes that the melted spots are antinodes, half a
  wavelength apart, and that the oven really runs at its labelled frequency.
  \textbf{Reliability:} repeat and average; measure across several gaps and divide.
  \textbf{Accuracy:} measure between spot centres. The spots are wide, so their centres are
  uncertain by a few mm, and the oven's frequency is only roughly \qty{2.45}{GHz}.}
\end{experiment}

\begin{experiment}{Spectra of discharge tubes, a filament lamp and sunlight}
\aim{to observe and compare the spectra of different light sources.}
\apparatus{spectroscope; H, He, Ne and Hg discharge tubes with a high-voltage supply; filament
  lamp; white card.
  \begin{figure}
    \centering
    \begin{tikzpicture}[x=1cm, y=1cm]
      \draw[line width=0.8pt, fill=hscray!35] (-0.85,0) circle (0.22);
      \node[font=\footnotesize, anchor=north] at (-0.85,-0.27) {source};
      \fill[hscink] (-0.06,0.1) rectangle (0.06,0.5);
      \fill[hscink] (-0.06,-0.1) rectangle (0.06,-0.5);
      \node[font=\footnotesize, anchor=north] at (0.05,-0.55) {slit};
      \draw[line width=0.6pt, hscmuted] (0.1,0.45) -- (2.3,0.45);
      \draw[line width=0.6pt, hscmuted] (0.1,-0.45) -- (2.3,-0.45);
      \draw[line width=0.8pt, fill=hscvel!15] (2.3,0) ellipse (0.08 and 0.43);
      \draw[ray] (0.06,0) -- (2.24,0.3);
      \draw[ray] (0.06,0) -- (2.24,-0.3);
      \draw[ray] (2.34,0.3) -- (3.75,0.3);
      \draw[ray] (2.34,-0.3) -- (3.75,-0.3);
      \node[font=\footnotesize, align=center, anchor=south] at (1.1,0.47)
        {collimator:\\makes the rays parallel};
      \draw[guide] (3.8,0) circle (0.95);
      \draw[line width=2pt, hscink] (3.8,-0.5) -- (3.8,0.5);
      \node[font=\footnotesize, anchor=north] at (3.8,-0.97) {grating on a turntable};
      \draw[guide] (3.8,0) -- (7.2,0) node[right, font=\footnotesize, align=left] {straight\\through};
      \begin{scope}[shift={(3.8,0)}, rotate=20]
        \draw[guide] (0,0) -- (1.2,0);
        \draw[ray] (0.05,0.28) -- (1.2,0.28);
        \draw[ray] (0.05,-0.28) -- (1.2,-0.28);
        \draw[line width=0.6pt, hscmuted] (1.25,0.42) -- (2.9,0.42);
        \draw[line width=0.6pt, hscmuted] (1.25,-0.42) -- (2.9,-0.42);
        \draw[line width=0.8pt, fill=hscvel!15] (1.25,0) ellipse (0.07 and 0.4);
        \draw[line width=0.8pt, fill=hscvel!15] (2.9,0) ellipse (0.05 and 0.24);
        \node[font=\footnotesize, anchor=south, rotate=20] at (2.05,0.45) {telescope};
      \end{scope}
      \draw[line width=0.6pt] (4.3,0) arc (0:20:0.5);
      \node[font=\scriptsize, inner sep=1pt] at ($(3.8,0)+(10:0.72)$) {$\theta$};
    \end{tikzpicture}
    \caption{A grating spectroscope. Each wavelength leaves the grating at its own angle
      $\theta$.}\label{ch1:fig:spectroscope}
  \end{figure}}
\method{in a dark room, view each discharge tube through the slit
  (\cref{ch1:fig:spectroscope}). Read the wavelength of each line from the spectroscope's scale
  (or measure each angle $\theta$ and use $d\sin\theta = m\lambda$, \cref{ch2:chap}). Then view
  the filament lamp, and sunlight reflected from a white card.}
\results{{\centering\small
  \begin{tabular}{@{}L{32mm}L{98mm}@{}}
    \toprule
    \colhead{Source} & \colhead{Spectrum seen}\\
    \midrule
    discharge tubes & emission: bright lines on black, a different set for each gas\\
    filament lamp & continuous: all colours, brightest at the red end\\
    reflected sunlight & continuous, with narrow dark lines, e.g.\ \qty{589}{nm} (Na), \qty{656}{nm} (H)\\
    \bottomrule
  \end{tabular}\par}}
\conclusion{hot dense sources give continuous spectra. Excited low-pressure gases give line
  spectra, unique to each element.}
\significance{the basis of spectroscopy in the laboratory and in astronomy.}
\validity{\textbf{Validity:} same spectroscope for every source; a dark room stops stray light.
  \textbf{Reliability:} repeat readings; compare with reference charts.
  \textbf{Accuracy:} calibrate the scale with known lines (mercury 436 and \qty{546}{nm}); use a
  narrow slit. \textbf{Safety:} several thousand volts, so switch off and let a tube cool before
  changing it. Mercury tubes give out UV. Never look at the Sun.}
\end{experiment}
